\section{Kth Missing Positive Number}
\label{sec:bs-kth-missing}

\problemstatement{We are given a sorted array $\textit{arr}$ of $n$
distinct positive integers, and an integer $k$. We must find the $k$-th
positive integer that is \textbf{missing} from $\textit{arr}$.}

\begin{patternbox}
There is no array index to ``find $X$ at'' here -- we are counting
\emph{missing} numbers, which is a different flavor of monotonic quantity
than anything in Sections~\ref{sec:bs-find-x}--\ref{sec:bs-peak-1d}.
Whenever a sorted-array problem asks about \emph{missing} elements (gaps in
an otherwise-expected sequence) and wants the $k$-th such gap, look for a
formula that counts ``how many integers are missing up to this point in
the array,'' and binary search for the boundary where that count first
reaches $k$.
\end{patternbox}

\subsection*{Understanding the problem carefully}

Consider index $i$ (0-indexed) in the array. If none of $\textit{arr}[0],
\dots, \textit{arr}[i]$ were missing from the positive integers
$1, \dots, \textit{arr}[i]$, there would be exactly $\textit{arr}[i]$ of
them present, occupying $\textit{arr}[i]$ slots; but only $i+1$ of them
\emph{are} actually present (the first $i+1$ elements of the array), so
\[
  \textit{missingCount}(i) = \textit{arr}[i] - (i+1)
\]
counts exactly how many positive integers from $1$ to $\textit{arr}[i]$ are
missing from the array. Because $\textit{arr}$ is sorted and strictly
increasing, $\textit{arr}[i]$ grows at least as fast as $i$ does, so
$\textit{missingCount}(i)$ is non-decreasing in $i$ -- exactly the
monotonic structure we need.

\begin{ideabox}[Candidate Space and Predicate]
The candidate space is indices $0, \dots, n-1$ (with an implicit
``beyond the array'' fallback). The predicate is
$P(i) = (\textit{missingCount}(i) \ge k)$, and we want the smallest index
where $P$ holds.
\end{ideabox}

\subsection*{Why $k + \textit{lo}$ is the final answer}

Let $\textit{lo}$ be the smallest index where $\textit{missingCount}
(\textit{lo}) \ge k$ (or $n$, if even the last element's missing count is
still below $k$, meaning the answer lies beyond the array entirely). Just
before this index, $\textit{missingCount}(\textit{lo}-1) < k$, meaning
fewer than $k$ numbers are missing among $1, \dots, \textit{arr}
[\textit{lo}-1]$; from $\textit{arr}[\textit{lo}-1]+1$ onward (which is not
yet present in the array, since $\textit{arr}[\textit{lo}]$ is strictly
larger), missing numbers resume counting up consecutively until reaching
$\textit{missingCount}(\textit{lo}) \ge k$. A clean way to see the final
formula: exactly $\textit{lo}$ array elements (indices $0, \dots,
\textit{lo}-1$) are known to be \emph{smaller} than the answer (since the
crossover into ``enough missing numbers'' has not yet happened at index
$\textit{lo}-1$), so among the first $k + \textit{lo}$ positive integers,
exactly $\textit{lo}$ are present in the array and the remaining $k$ are
missing -- meaning the $k$-th missing positive integer is precisely
$k + \textit{lo}$.

\subsection*{Algorithm}

\begin{enumerate}
  \item Initialize $\textit{lo} \gets 0$, $\textit{hi} \gets n-1$.
  \item While $\textit{lo} \le \textit{hi}$:
  \begin{enumerate}
    \item $\textit{mid} \gets \textit{lo}+(\textit{hi}-\textit{lo})/2$.
    \item Compute $\textit{missingCount} \gets \textit{arr}[\textit{mid}] -
          (\textit{mid}+1)$.
    \item If $\textit{missingCount} < k$: $\textit{lo} \gets
          \textit{mid}+1$.
    \item Else: $\textit{hi} \gets \textit{mid}-1$.
  \end{enumerate}
  \item Return $k + \textit{lo}$.
\end{enumerate}

\begin{algorithm}[H]
\caption{Kth Missing Positive Number}
\begin{algorithmic}[1]
\Procedure{KthMissing}{$\textit{arr}, k$}
  \State $\textit{lo} \gets 0$, $\textit{hi} \gets n-1$
  \While{$\textit{lo} \le \textit{hi}$}
    \State $\textit{mid} \gets \textit{lo}+(\textit{hi}-\textit{lo})/2$
    \State $\textit{missing} \gets \textit{arr}[\textit{mid}] - (\textit{mid}+1)$
    \If{$\textit{missing} < k$}
      \State $\textit{lo} \gets \textit{mid} + 1$
    \Else
      \State $\textit{hi} \gets \textit{mid} - 1$
    \EndIf
  \EndWhile
  \State \Return $k + \textit{lo}$
\EndProcedure
\end{algorithmic}
\end{algorithm}

\subsection*{Dry run}

\dryrun{Take $\textit{arr} = [2,3,4,7,11]$, $k = 5$.}

\begin{itemize}
  \item $\textit{lo}=0,\textit{hi}=4$: $\textit{mid}=2$,
        $\textit{missing}=\textit{arr}[2]-3=4-3=1 < 5$. $\textit{lo}=3$.
  \item $\textit{lo}=3,\textit{hi}=4$: $\textit{mid}=3$,
        $\textit{missing}=7-4=3 < 5$. $\textit{lo}=4$.
  \item $\textit{lo}=4,\textit{hi}=4$: $\textit{mid}=4$,
        $\textit{missing}=11-5=6 \ge 5$. $\textit{hi}=3$.
  \item $\textit{lo}=4 > \textit{hi}=3$: loop ends.
\end{itemize}

The answer is $k+\textit{lo} = 5+4 = 9$. Checking by hand: the missing
positive integers, in order, are $1, 5, 6, 8, 9, 10, 12, \dots$, and the
$5$-th one is indeed $9$.

\subsection*{C++ implementation}

\begin{lstlisting}
#include <bits/stdc++.h>
using namespace std;

int findKthMissing(vector<int>& arr, int k) {
    int n = arr.size();
    int lo = 0, hi = n - 1;

    while (lo <= hi) {
        int mid = lo + (hi - lo) / 2;
        int missing = arr[mid] - (mid + 1);
        if (missing < k) {
            lo = mid + 1;
        } else {
            hi = mid - 1;
        }
    }
    return k + lo;
}

int main() {
    vector<int> arr = {2, 3, 4, 7, 11};
    cout << "5th missing positive number: "
         << findKthMissing(arr, 5) << endl;
    return 0;
}
\end{lstlisting}

\begin{complexitybox}
\textbf{Time Complexity.} The search window halves each iteration, giving
\[
  O(\log n).
\]
This is a substantial improvement over the natural $O(n)$ linear-scan
approach of walking through the array while counting missing numbers one
at a time.
\textbf{Space Complexity.} $O(1)$ auxiliary space.
\end{complexitybox}

\begin{takeawaybox}
When a problem asks about ``missing'' elements relative to an expected
dense sequence (here, the positive integers), look for a closed-form
counting function -- ``expected count minus actual count up to this
point'' -- that is monotonic, and binary search on that function's
crossover rather than scanning and counting one gap at a time.
\end{takeawaybox}
